\subsection{选择校准规模}

为了分辨千分之一量级的参考不一致，\cref{sec:6-2} 的独立槽位模型提示
\[
 \frac{4.74}{Sk}\simeq10^{-3},\quad k=10
 \quad\Longrightarrow\quad S\simeq474\approx500.
\]
这为除小样本外测试整个 500 条查询的校准池提供了理由。该模型针对参考一致性；实测留出质量及抽样风险见 \cref{tab:7-budget-full,fig:7-budget-risk}。全池二分的搜索顺序条件仍为 \cref{sec:6-3} 所述条件。

\subsection{风险与所选预算性能}

令 $m$ 个随机校准样本分别选择预算 $\widehat\alpha_r$。以最佳已测留出召回率为参考，定义第 $r$ 个样本的损失及损失超过目标的比例：
\begin{equation}
 L_r=\widehat r^*-\widehat r(\widehat\alpha_r),\qquad
 \widehat p_\delta=\frac{1}{m}\sum_{r=1}^{m}\mathbf1\{L_r>\delta\}.
 \label{eq:6-risk}
\end{equation}
平均损失说明平均质量，越阈频率说明多常未达标，尾部损失则说明未达标时的幅度。\Cref{tab:7-budget-quality,fig:7-budget-risk} 报告这些汇总，逐样本最大损失见 \cref{sec:7-budget-trace}。

报告速率 $\overline Q=m^{-1}\sum_r\widehat Q(\widehat\alpha_r)$ 将各样本选中预算的既有实测吞吐取平均。这是基于固定预算运行的估计，不同于中位预算的速率，也不同于直接计时校准后的查询流。\Cref{tab:7-budget-cost} 将其与质量一起报告。

\subsection{复用成本}

只有后续搜索节省足够时间，校准才能回本。令 $T_{\rm cal}$ 为一次性校准成本，$T_{\rm chosen}$ 为所选预算下每 batch 耗时，$T_f$ 为固定预算下同样 batch 的耗时。对后续 $B_{\rm reuse}$ 个等大小 batch，
\begin{equation}
 T_{\rm policy}(B_{\rm reuse})=T_{\rm cal}+B_{\rm reuse}T_{\rm chosen},\qquad
 T_{\rm fixed}(B_{\rm reuse})=B_{\rm reuse}T_f.
 \label{eq:6-total-cost}
\end{equation}
当 $T_{\rm chosen}<T_f$ 时，用启动成本除以每 batch 节省时间，得到回本次数：
\begin{equation}
 B_{\rm reuse}^*=\frac{T_{\rm cal}}{T_f-T_{\rm chosen}}.
 \label{eq:6-payback}
\end{equation}
严格节省时间的条件直接由下式得到：
\[
 \begin{aligned}
 T_{\rm policy}<T_{\rm fixed}
 &\iff T_{\rm cal}<B_{\rm reuse}(T_f-T_{\rm chosen})\\
 &\iff B_{\rm reuse}>B_{\rm reuse}^*
 \qquad(T_{\rm chosen}<T_f).
 \end{aligned}
\]
若 $T_{\rm chosen}\geq T_f$，则无论多少 batch 都不能靠吞吐收回校准成本。每 batch 有 $n$ 条查询时，$T=1000n/Q$ 将 QPS 转为毫秒。当所选预算变化时，$1/\mathbb E[Q]\leq\mathbb E[1/Q]$，取平均 QPS 的倒数会低估平均 batch 时间。计时边界和直接查询流检查见 \cref{sec:a-budget-contract,sec:7-stream-cost}，模型估计见 \cref{tab:7-budget-cost}。

每隔 $H$ 个 batch 重新校准，会增加平均每 batch 的 $T_{\rm cal}/H$ 成本。自动检测工作负载变化及更新校准仍属后续工作。
